\subsection{Undoing differentiation and measuring total change}

\begin{corebox}[Two connected meanings]
An indefinite integral asks for antiderivatives. A definite integral measures accumulated signed change over an interval. The fundamental theorem will connect these two meanings.
\end{corebox}

\begin{definitionbox}[Antiderivative]
A function $F$ is an antiderivative of $f$ when
\[
F'(x)=f(x).
\]
We write
\[
\int f(x)\,dx=F(x)+C.
\]
\end{definitionbox}

The constant $C$ is necessary because differentiation removes additive constants.

\begin{definitionbox}[Definite integral]
The expression
\[
\int_a^b f(x)\,dx
\]
represents accumulated signed contribution of $f$ from $a$ to $b$.
\end{definitionbox}

\begin{examplebox}[The same integrand, two questions]
\[
\int 2x\,dx=x^2+C,
\]
while
\[
\int_0^3 2x\,dx=9.
\]
The first answer is a family of functions; the second is a number.
\end{examplebox}

\begin{interpretationbox}[Accumulation from small pieces]
The integral combines many small contributions. For a non-negative function this can be seen as area; for a rate it represents total change.
\end{interpretationbox}
